\chapter{Увод}

\section{Описание и идея на проекта}
Основната идея на проекта е създаването на конзолно приложение на езика \textit{C++} за управление на електронни таблици. Програмата предоставя текстов интерфейс, чрез който потребителят може да зарежда данни от файл в паметта, да редактира съдържанието на отделните клетки и да визуализира таблицата. Поддържа се работа с различни типове данни в една и съща таблица, както и динамично изчисляване на математически формули с препратки към други клетки. Всички промени се запазват обратно в текстов файл при изрично желание на потребителя.

Проектът е с отворен код и достъпен в GitHub на адрес: \\ \url{https://github.com/AVachkov/Spreadsheet}.

\section{Цел и задачи на разработката}

Целта на проекта е да се изгради работеща система, която абстрахира сложността на вътрешното представяне на данните и се фокусира изцяло върху същността на тяхната обработка, изчисляване и сигурно съхранение.

Основните задачи, решени по време на разработката, са:

\begin{itemize}
  \item \textbf{Структура на данните:} Създаване на общ базов клас и отделни класове за поддържаните хетерогенни типове съдържание: цели числа, дробни числа, символни низове и формули.
  \item \textbf{Взаимодействие между класовете:} Правилно организиране на връзките и комуникацията между главния клас на таблицата и отделните обекти за клетките, за да работят заедно като единна система.
  \item \textbf{Пресмятане на формули:} Реализиране на алгоритъм, който чете математически изрази, поддържа основните аритметични операции (\texttt{+}, \texttt{-}, \texttt{*}, \texttt{/}, \texttt{\^{}}), спазва техния приоритет и преобразува типовете при нужда (например празни клетки и нечислови низове се оценяват като 0).
  \item \textbf{Изчисляване на препратките:} Извличане на стойности от конкретни адреси (например \texttt{R1C1}), като при грешка във формулата или деление на нула програмата не прекъсва своето изпълнение, а просто извежда \texttt{"ERROR"} в съответната клетка.
  \item \textbf{Защита от зацикляне:} Справяне с проблема, при който две клетки зависят една от друга, за да се предотврати блокиране на програмата.
  \item \textbf{Управление на паметта:} Внимателно и оптимално заделяне и освобождаване на динамична памет, за да няма утечки (\textit{memory leak}).
  \item \textbf{Работа с файлове:} Реализиране на команди за зареждане (\texttt{open}) и записване (\texttt{save}, \texttt{saveas}) на таблицата.
\end{itemize}

\section{Структура на документацията}
Настоящата документация е организирана в пет основни глави:
\begin{itemize}
  \item \textbf{Глава 1 - Увод:} Описва идеята на проекта, целите и основните решени задачи.
  \item \textbf{Глава 2 - Преглед на предметната област:} Разглежда дефинициите за типовете данни, формулите и функционалните изисквания.
  \item \textbf{Глава 3 - Проектиране:} Представя архитектурата на приложението и ООП дизайна.
  \item \textbf{Глава 4 - Реализация и тестване:} Описва детайли по управлението на паметта, алгоритмите и проведените тестове.
  \item \textbf{Глава 5 - Заключение:} Обобщава постигнатите резултати и дава насоки за бъдещо развитие.
\end{itemize}